\documentclass[12pt,a4paper]{article}

% --- Packages -----------------------------------------------------------
\usepackage[utf8]{inputenc}
\usepackage[T1]{fontenc}
\usepackage[a4paper, margin=2.5cm]{geometry}
\usepackage{lmodern}
\usepackage{microtype}
\usepackage{parskip}
\usepackage{titlesec}
\usepackage{titletoc}
\usepackage{tocloft}
\usepackage{hyperref}
\usepackage{xcolor}
\usepackage{booktabs}
\usepackage{longtable}
\usepackage{tabularx}
\usepackage{array}
\usepackage{colortbl}
\usepackage{enumitem}
\usepackage{fancyhdr}
\usepackage{graphicx}
\usepackage{mdframed}
\usepackage{listings}
\usepackage{amsmath}
\usepackage{multirow}
\usepackage{makecell}
\usepackage{float}
\usepackage{caption}
\usepackage{subcaption}
\usepackage{soul}

% --- Colours ------------------------------------------------------------
\definecolor{zaraiGreen}{HTML}{059669}
\definecolor{zaraiBlue}{HTML}{4F46E5}
\definecolor{zaraiRed}{HTML}{DC2626}
\definecolor{zaraiAmber}{HTML}{D97706}
\definecolor{zaraiGray}{HTML}{374151}
\definecolor{codeBg}{HTML}{F3F4F6}
\definecolor{warningBg}{HTML}{FFFBEB}
\definecolor{warningBorder}{HTML}{F59E0B}
\definecolor{errorBg}{HTML}{FEF2F2}
\definecolor{errorBorder}{HTML}{EF4444}
\definecolor{okBg}{HTML}{ECFDF5}
\definecolor{okBorder}{HTML}{10B981}

% --- Hyperref setup ------------------------------------------------------
\hypersetup{
  colorlinks=true,
  linkcolor=zaraiBlue,
  urlcolor=zaraiGreen,
  citecolor=zaraiGray,
  pdfauthor={Zarailink FYP Team},
  pdftitle={Zarailink Search Engine: Exhaustive Technical Assessment},
}

% --- Section formatting --------------------------------------------------
\titleformat{\section}{\large\bfseries\color{zaraiGray}}{}{0em}{\thesection\quad}[\titlerule]
\titleformat{\subsection}{\normalsize\bfseries\color{zaraiBlue}}{}{0em}{\thesubsection\quad}
\titleformat{\subsubsection}{\normalsize\bfseries\color{zaraiGreen}}{}{0em}{\thesubsubsection\quad}
\titlespacing*{\section}{0pt}{1.5em}{0.6em}
\titlespacing*{\subsection}{0pt}{1em}{0.4em}
\titlespacing*{\subsubsection}{0pt}{0.8em}{0.3em}

% --- Header / Footer -----------------------------------------------------
\pagestyle{fancy}
\fancyhf{}
\fancyhead[L]{\small\color{zaraiGray}\textit{Zarailink Search Engine --- Exhaustive Assessment}}
\fancyhead[R]{\small\color{zaraiGray}\thepage}
\fancyfoot[C]{\small\color{zaraiGray}Confidential --- FYP Internal Report}
\renewcommand{\headrulewidth}{0.4pt}
\renewcommand{\footrulewidth}{0.4pt}

% --- Code listings -------------------------------------------------------
\lstset{
  backgroundcolor=\color{codeBg},
  basicstyle=\small\ttfamily\color{zaraiGray},
  keywordstyle=\bfseries\color{zaraiBlue},
  commentstyle=\itshape\color{gray},
  stringstyle=\color{zaraiGreen},
  numbers=none,
  frame=none,
  breaklines=true,
  breakatwhitespace=true,
  tabsize=4,
  showstringspaces=false,
  xleftmargin=1em,
  xrightmargin=1em,
}

% --- Custom environments -------------------------------------------------
\newmdenv[
  backgroundcolor=errorBg,
  linecolor=errorBorder,
  linewidth=1.5pt,
  leftmargin=0pt,
  rightmargin=0pt,
  innerleftmargin=12pt,
  innerrightmargin=12pt,
  innertopmargin=8pt,
  innerbottommargin=8pt,
  skipabove=6pt,
  skipbelow=6pt,
]{issueBox}

\newmdenv[
  backgroundcolor=warningBg,
  linecolor=warningBorder,
  linewidth=1.5pt,
  leftmargin=0pt,
  rightmargin=0pt,
  innerleftmargin=12pt,
  innerrightmargin=12pt,
  innertopmargin=8pt,
  innerbottommargin=8pt,
  skipabove=6pt,
  skipbelow=6pt,
]{warnBox}

\newmdenv[
  backgroundcolor=okBg,
  linecolor=okBorder,
  linewidth=1.5pt,
  leftmargin=0pt,
  rightmargin=0pt,
  innerleftmargin=12pt,
  innerrightmargin=12pt,
  innertopmargin=8pt,
  innerbottommargin=8pt,
  skipabove=6pt,
  skipbelow=6pt,
]{okBox}

% --- Table cell colours ---------------------------------------------------
\newcommand{\cellOK}[1]{\cellcolor{okBg}\textcolor{zaraiGreen}{\textbf{#1}}}
\newcommand{\cellWarn}[1]{\cellcolor{warningBg}\textcolor{zaraiAmber}{\textbf{#1}}}
\newcommand{\cellBad}[1]{\cellcolor{errorBg}\textcolor{zaraiRed}{\textbf{#1}}}

% ------------------------------------------------------------------------
\begin{document}

% --- Title page ----------------------------------------------------------
\begin{titlepage}
  \centering
  \vspace*{2cm}
  {\Huge\bfseries\color{zaraiGray} Zarailink Search Engine}\\[0.5em]
  {\Large\bfseries\color{zaraiBlue} Exhaustive Technical Assessment Report}\\[2em]
  \rule{10cm}{0.8pt}\\[1.5em]
  {\large\color{zaraiGray}
    Covers every layer of the search pipeline: \\
    NLU \textbullet{}  Product Resolution \textbullet{}  Aggregation \textbullet{}  Ranking \textbullet{}  Frontend \textbullet{}  \\
    Worldwide/Pakistan Scope \textbullet{}  Edge Cases \textbullet{}  Bugs \textbullet{}  Gaps}\\[2em]
  \rule{10cm}{0.4pt}\\[2em]
  {\normalsize\color{gray}
    Final Year Project --- 8th Semester\\
    Zarailink: Pakistani Agricultural Trade Intelligence Platform\\[0.5em]
    \textbf{Date:} April 17, 2026\\
    \textbf{Classification:} Internal / Confidential
  }
  \vfill
  \includegraphics[width=2.5cm]{example-image}\\[0.5em]
  {\small\color{gray}\textit{(replace placeholder with Zarailink logo)}}
\end{titlepage}

% --- Table of Contents ---------------------------------------------------
\tableofcontents
\newpage

% ========================================================================
\section{Executive Summary}
% ========================================================================

The Zarailink search engine is a four-stage NLP pipeline built for B2B agricultural trade
intelligence on Pakistani import/export data. After exhaustive review of every file in the
search module --- \texttt{nlu\_engine.py} (941 lines), \texttt{search\_service.py} (790 lines),
\texttt{aggregation.py} (625 lines), \texttt{views.py} (260 lines), and both frontend
components --- the following verdict emerges:

\begin{table}[H]
\centering
\renewcommand{\arraystretch}{1.4}
\begin{tabularx}{\textwidth}{lXl}
\toprule
\textbf{Layer} & \textbf{Summary} & \textbf{Grade} \\
\midrule
Intent Detection & Works for explicit verbs; breaks on implicit phrasing and Urdu &\cellWarn{B-} \\
Country Extraction & Handles ~20 countries; many countries silently undetected &\cellWarn{C+} \\
Product Keyword Extraction & Stop-word strip works for simple queries; multi-word qualifiers break it &\cellWarn{C+} \\
Product Resolution (Fuzzy Matching) & 4-pass system is solid; BUT full-keyword search before word decomposition hurts &\cellWarn{B} \\
Scope / Worldwide/Pakistan & Correct architecture; UI metaphor is confusing without explanation &\cellWarn{B} \\
Aggregation Query & Excellent --- correct semantics, proper weighted avg price &\cellOK{A-} \\
Ranking Model & Heuristic only (LTR model deleted); weights are simplistic &\cellBad{C} \\
Caching & Three-level cache is well-designed &\cellOK{A-} \\
Frontend UI & Clean, functional; several state and filter logic bugs &\cellWarn{B-} \\
Overall Search Quality & Handles 60-70\% of realistic queries correctly end-to-end &\cellWarn{C+} \\
\bottomrule
\end{tabularx}
\caption{Layer-by-layer assessment grades}
\end{table}

\textbf{The single largest problem:} The LightGBM LTR ranking model referenced in the
project documentation was \textit{deleted} (\texttt{ranking.py}, \texttt{ranking\_ltr.py}
removed in a prior refactor). The active ranking is a purely heuristic min-max weighted
composite with no learned signal whatsoever.

\textbf{Second largest problem:} The NLU pipeline has critical blind spots on any query
that doesn't use explicit English trade verbs --- non-English queries, casual phrasing,
product qualifiers, company name searches, and HS code searches all fail silently.

% ========================================================================
\section{System Architecture Overview}
% ========================================================================

The search engine follows a strictly sequential pipeline. Every request through the
\texttt{GET /api/search/?q=\ldots} endpoint passes through these stages:

\begin{center}
\begin{tabular}{ccccc}
\textbf{Stage 1} & $\rightarrow$ & \textbf{Stage 2} & $\rightarrow$ & \textbf{Stage 3} \\
NLU Parse & & Product Resolution & & ORM Aggregation \\
\small(intent+keyword+country+price) & & \small(4-pass fuzzy match) & & \small(Django QuerySet) \\
\end{tabular}

\vspace{0.5em}
$\downarrow$

\begin{tabular}{c}
\textbf{Stage 4} \\
Ranking + Response \\
\small(heuristic composite score)
\end{tabular}
\end{center}

\vspace{1em}
\textbf{Key design decisions:}
\begin{itemize}
  \item OpenSearch/BM25 path exists in code but is permanently \texttt{if False} disabled.
  \item Three-level caching: NLU (24h), availability (1h), aggregation (1h).
  \item Scope (Worldwide/Pakistan) is passed as a UI parameter that determines which
        trade type (IMPORT/EXPORT) to query --- it is not inferred from the query itself.
  \item The LTR model files exist at \texttt{search/models/lgbm\_ltr.txt} but the code
        that loads and uses them (\texttt{ranking.py}, \texttt{ranking\_ltr.py}) was deleted.
\end{itemize}

% ========================================================================
\section{Stage 1: NLU Engine --- Deep Analysis}
% ========================================================================

\subsection{Intent Detection}

\subsubsection{What Works Correctly}

Intent detection follows a four-level priority cascade defined in \texttt{nlu\_engine.py:549}:

\begin{enumerate}
  \item \textbf{Pattern rules} --- semantic reversal detection (e.g.\ \textit{"find sellers of X"} $\rightarrow$ BUY)
  \item \textbf{SetFit ML model} --- if trained model exists at \texttt{search/models/intent\_model/}
  \item \textbf{Regex keyword fallback} --- scans for explicit trade verbs
  \item \textbf{Default: BUY}
\end{enumerate}

The pattern rules correctly handle the most confusing class of queries:

\begin{lstlisting}[language=Python]
# "looking for sellers of rice" $\rightarrow$ BUY (user wants to buy from those sellers)
if re.search(r'\b(?:find|need|want|looking for)\s+(?:sellers?|suppliers?)\b', q):
    return "BUY"

# "find buyers of wheat" $\rightarrow$ SELL (user wants to sell to those buyers)
if re.search(r'\b(?:find|need|want|looking for)\s+(?:buyers?|importers?)\b', q):
    return "SELL"
\end{lstlisting}

The regex fallback correctly classifies explicit-verb queries: \textit{"i want to buy sugar"} $\rightarrow$ BUY, \textit{"sell wheat"} $\rightarrow$ SELL.

\subsubsection{Intent Detection: Critical Failures}

\begin{issueBox}
\textbf{CRITICAL: Urdu and non-English queries are completely unhandled.}
The intent patterns, stop-word lists, and regex are 100\% English. A user typing
\textit{``cheeni kharidna''} (Urdu transliteration: buy sugar) gets intent=BUY by default
--- but only by luck. No character or token is understood; the Urdu script query is
noise to the entire pipeline. Arabic-script Urdu queries return 0 results every time.
\end{issueBox}

\begin{issueBox}
\textbf{ISSUE: ``for sale'' phrasing returns wrong intent.}
Query: \textit{"dextrose for sale"} \\
Expected intent: SELL (user has product for sale). \\
Actual: BUY (no sell verb matched; default is BUY). \\
\texttt{"sale"} does not match the regex \texttt{sell\textbackslash w*} because ``sale'' $\neq$ ``sell+suffix''.
\end{issueBox}

\begin{issueBox}
\textbf{ISSUE: Noun-phrase selling signals ignored.}
Queries like \textit{"dextrose available"}, \textit{"sugar stock to offload"},
\textit{"wheat inventory"}, \textit{"we have PVC resin"} all return BUY by default.
None of these trigger the SELL regex.
\end{issueBox}

\begin{issueBox}
\textbf{ISSUE: ``I am a supplier of X'' returns BUY.}
The regex \texttt{supplier\textbackslash w*} maps to BUY (you are looking for a supplier).
But \textit{"I am a supplier of rice"} means the user wants to SELL. The regex has no
way to distinguish ``I am a supplier'' from ``I need a supplier.''
\end{issueBox}

\begin{warnBox}
\textbf{WARNING: SetFit model presence is unverified at runtime.}
The code silently falls back to regex if the model is missing. There is no health-check
endpoint or admin indicator telling operators whether the ML model is active or not.
\end{warnBox}

\subsubsection{Intent Detection --- Query Coverage Matrix}

\begin{table}[H]
\centering
\renewcommand{\arraystretch}{1.3}
\begin{tabularx}{\textwidth}{Xcc}
\toprule
\textbf{Query Example} & \textbf{Expected} & \textbf{Actual} \\
\midrule
\textit{buy dextrose} & BUY & \cellOK{BUY} \\
\textit{sell sugar} & SELL & \cellOK{SELL} \\
\textit{find sellers of wheat} & BUY & \cellOK{BUY} \\
\textit{find buyers for rice} & SELL & \cellOK{SELL} \\
\textit{import urea from china} & BUY & \cellOK{BUY} \\
\textit{export wheat to UAE} & SELL & \cellOK{SELL} \\
\textit{dextrose anhydrous} (no verb) & BUY (default ok) & \cellOK{BUY} \\
\textit{who sells PVC resin} & BUY & \cellOK{BUY} \\
\textit{dextrose for sale} & SELL & \cellBad{BUY} \\
\textit{I am a supplier of rice} & SELL & \cellBad{BUY} \\
\textit{sugar available to export} & SELL & \cellOK{SELL} \\
\textit{we have wheat in stock} & SELL & \cellBad{BUY} \\
\textit{dextrose supplier needed} & BUY & \cellOK{BUY} \\
\textit{Urdu-script query: ``buy sugar''} & BUY & \cellWarn{BUY (luck)} \\
\textit{1702.3090} (HS code only) & BUY & \cellWarn{BUY (luck)} \\
\textit{SEAWALL} (company name) & N/A & \cellBad{no results} \\
\bottomrule
\end{tabularx}
\caption{Intent detection accuracy across query types}
\end{table}

\subsection{Product Keyword Extraction}

\subsubsection{The Stop-Word Stripping Mechanism}

The primary extraction method in \texttt{extract\_product\_keyword()} (line 84) performs:
\begin{enumerate}
  \item Regex strip of garbled intent verbs: \texttt{buy\textbackslash w*, sell\textbackslash w*, get\textbackslash w*, import\textbackslash w*\ldots}
  \item Sequential removal of 40+ stop phrases from a hardcoded list
  \item Price clause removal: \texttt{under \$700}, \texttt{above 500 usd}
  \item Tokenization fallback for long leftover strings (trigram DB scoring)
\end{enumerate}

\subsubsection{What Works}

\begin{okBox}
Good: \textit{"i want to buy dextrose anhydrous"} $\rightarrow$ \textit{"dextrose anhydrous"} \\
Good: \textit{"find sellers of basmati rice"} $\rightarrow$ \textit{"basmati rice"} \\
Good: \textit{"sugar under \$500"} $\rightarrow$ \textit{"sugar"} \\
Good: \textit{"buyyy deeeex cheap"} $\rightarrow$ regex strips garbled verbs $\rightarrow$ \textit{"deeeex cheap"} $\rightarrow$ trigram finds \textit{"Dextrose"}
\end{okBox}

\subsubsection{Product Keyword Extraction: Critical Failures}

\begin{issueBox}
\textbf{CRITICAL: Business qualifiers before product name break extraction.}
\textit{"wholesale sugar"} $\rightarrow$ \textit{"wholesale sugar"} (not stripped; ``wholesale'' not in stop words). \\
\textit{"organic wheat"} $\rightarrow$ \textit{"organic wheat"} (not stripped). \\
These then fail prefix/contains DB match and must rely on trigram with keyword=\textit{"organic wheat"}.
Trigram for \textit{"organic wheat"} vs DB entry \textit{"Wheat"}: partial match at $\approx$0.4 similarity. Might work,
but is fragile --- one new product named \textit{"Organic Compound"} could outscore it.
\end{issueBox}

\begin{issueBox}
\textbf{CRITICAL: No support for HS code as the search query.}
If a user types \textit{"1702.3090"} or \textit{"17021990"}, the product keyword after
stripping is the HS code itself. No DB lookup tries to match HS codes from the query
string in the main NLU path. The HS code is only used when passed as a separate URL
parameter (e.g.\ selected from autocomplete). Users who know HS codes get 0 results.
\end{issueBox}

\begin{issueBox}
\textbf{ISSUE: Multi-product queries treated as a single keyword.}
\textit{"sugar and wheat"} $\rightarrow$ \textit{"sugar and wheat"} (``and'' not stripped). \\
The system cannot return results for two products simultaneously.
Trigram for \textit{"sugar and wheat"} might find \textit{"Sugar"} at ~0.38 similarity, or nothing.
\end{issueBox}

\begin{issueBox}
\textbf{ISSUE: Tokenization fallback picks wrong word for uncommon qualifiers.}
Query: \textit{"best grade refined sugar"} \\
After stripping ``best'': \textit{"grade refined sugar"} (3 words, not $>4$, so no tokenization fallback). \\
DB lookup for \textit{"grade refined sugar"} fails prefix/contains; trigram similarity to \textit{"Sugar"}
may pass but is unreliable.
\end{issueBox}

\begin{warnBox}
\textbf{WARNING: KeyBERT confidence gate is too aggressive.}
When KeyBERT confidence $>$ 0.85 with no numbers or price operators, the LLM call is
entirely skipped. This means for a query like \textit{"premium dextrose anhydrous bulk"},
KeyBERT extracts \textit{"dextrose anhydrous"} with high confidence. But \textit{"bulk"} is a
ranking signal (\texttt{ranking\_hint=price\_desc}) that the regex price extractor does NOT
catch by itself (the regex only detects number-based prices). The bulk ranking hint is
in the LLM prompt but the regex has no ``bulk'' handler. So the supplier list is ranked
by default weights, not bulk-preference weights, silently.
\end{warnBox}

\begin{issueBox}
\textbf{BUG in confidence gate for ``bulk'' and similar signals:}
In \texttt{nlu\_engine.py:728}:
\begin{lstlisting}[language=Python]
ranking_word_hit = any(w in query.lower() for w in {
    "cheap", "cheapest", "bulk", "premium", "affordable", ...
})
\end{lstlisting}
This variable is checked in the confidence gate condition. If \textit{"bulk"} is present,
the LLM IS called --- good. But the \textit{regex fallback} (\texttt{\_build\_price\_filter\_regex})
only handles price-number patterns and \textit{"cheap/affordable"} as ranking hints.
\textit{"bulk"} is not in the regex's no-number ranking word list:
\begin{lstlisting}[language=Python]
price_ranking_words = {"cheap", "affordable", "budget", "cheapest",
                       "low price", "best price", "inexpensive"}
# NOTE: "bulk", "premium", "reliable" are MISSING
\end{lstlisting}
So if the LLM fails (timeout/API key missing), \textit{"bulk dextrose"} gets no
\texttt{ranking\_hint} and uses default weights, silently degrading quality.
\end{issueBox}

\subsection{Country Extraction}

\subsubsection{The Three-Layer Approach}

Country extraction uses: (1) FALLBACK\_COUNTRIES dict (regex), (2) GLiNER NER,
(3) RapidFuzz fuzzy matching. The regex dict is checked first.

\subsubsection{Country Coverage Analysis}

\begin{table}[H]
\centering
\renewcommand{\arraystretch}{1.3}
\begin{tabularx}{\textwidth}{Xcc}
\toprule
\textbf{Query / Country} & \textbf{Expected} & \textbf{Extracted} \\
\midrule
\textit{sugar from china} & China & \cellOK{China} \\
\textit{pak sugar supplier} & Pakistan & \cellOK{Pakistan} \\
\textit{ksa import} & Saudi Arabia & \cellOK{Saudi Arabia} \\
\textit{turk wheat} & Turkey & \cellOK{Turkey} \\
\textit{uae dextrose} & UAE & \cellOK{UAE} \\
\textit{bangladesh wheat} & Bangladesh & \cellBad{None (not in list)} \\
\textit{myanmar rice} & Myanmar & \cellBad{None (not in list)} \\
\textit{philippines coconut} & Philippines & \cellBad{None (not in list)} \\
\textit{morocco phosphate} & Morocco & \cellBad{None (not in list)} \\
\textit{hong kong} & Hong Kong & \cellBad{None (not in list)} \\
\textit{taiwan sugar} & Taiwan & \cellBad{None (not in list)} \\
\textit{russia wheat} & Russia & \cellBad{None (not in list)} \\
\textit{netherlands chemicals} & Netherlands & \cellBad{None (not in list)} \\
\textit{iranian sugar} & Iran & \cellWarn{Fuzzy may fail (``iranian'')} \\
\textit{australia canola} & Australia & \cellOK{Australia (via RapidFuzz)} \\
\bottomrule
\end{tabularx}
\caption{Country extraction accuracy across query types (None = country filter silently dropped)}
\end{table}

\begin{issueBox}
\textbf{CRITICAL: STANDARD\_COUNTRIES list is only 27 countries.}
Bangladesh, Philippines, Myanmar, Morocco, Hong Kong, Taiwan, Russia, Netherlands,
Poland, Ukraine, Israel, Ethiopia, Ghana, and many more are completely undetected.
If a user searches \textit{"bangladesh wheat supplier"}, the query is treated as if
no country was specified --- country filter is silently dropped and ALL wheat suppliers
appear, defeating the user's intent.
\end{issueBox}

\begin{warnBox}
\textbf{WARNING: RapidFuzz score\_cutoff asymmetry.}
The word-level RapidFuzz scan in \texttt{\_detect\_country\_fallback} uses score\_cutoff=85.0
while the GLiNER-extracted string resolver uses 60.0. This means a direct word like
\textit{"germany"} must score $\geq$85 to be accepted (it will), but
\textit{"germani"} (typo) scores $\approx$78 and is rejected at the word level.
The GLiNER path would catch it at 60.0, but GLiNER may not extract it as a ``country''
entity at all for a misspelling.
\end{warnBox}

\subsubsection{Country Extraction Edge Cases}

\begin{warnBox}
\textbf{False-positive risk:} \textit{"jordan almonds"} --- \textit{"jordan"} has no entry in
FALLBACK\_COUNTRIES but passes as word to RapidFuzz. WRatio("jordan", "South Korea")$\approx$40,
("jordan", "Jordan")? Jordan is NOT in STANDARD\_COUNTRIES. So it would correctly fail.
But \textit{"iran"} IS in FALLBACK\_COUNTRIES as ``Iran'' while \textit{"Iranian"} is not.
The word-level check would test \textit{"iranian"} against RapidFuzz and score against
\textit{"India"} at WRatio$\approx$55 (below 85 cutoff) --- so it correctly fails. Good.
But if GLiNER extracts \textit{"Iranian"} as a country and it goes to \_resolve\_country,
fuzzy against \textit{"India"} at 60 cutoff... WRatio("iranian","india")$\approx$47, OK.
Against \textit{"Iran"}: Iran is NOT in STANDARD\_COUNTRIES list either. Iran is missing entirely.
\end{warnBox}

\begin{issueBox}
\textbf{BUG: Iran is in FALLBACK\_COUNTRIES as a key but MISSING from STANDARD\_COUNTRIES.}
This means: exact word \textit{"iran"} $\rightarrow$ detected via fallback dict $\rightarrow$ ``Iran''. \\
But \textit{"iranian"} or GLiNER-extracted \textit{"Iran"} $\rightarrow$ goes to \_resolve\_country $\rightarrow$
no fuzzy match (Iran not in STANDARD\_COUNTRIES) $\rightarrow$ returns \textit{"Iran".capitalize()}
= \textit{"Iran"} via last-resort fallback. Accidentally correct but for wrong reasons.
\end{issueBox}

\subsection{Price Extraction}

\subsubsection{What Works}

The regex price extractor correctly handles:
\begin{itemize}
  \item \textit{"under \$700"}, \textit{"below 500 usd"}, \textit{"cheaper than 800"} $\rightarrow$ ceiling filter
  \item \textit{"above 1000"}, \textit{"at least 500"} $\rightarrow$ floor filter
  \item \textit{"around \$500"} $\rightarrow$ $\pm$15\% range (425--575)
  \item \textit{"cheap"}, \textit{"affordable"} (no number) $\rightarrow$ ranking\_hint only
\end{itemize}

\subsubsection{Price Extraction Failures}

\begin{issueBox}
\textbf{ISSUE: Written-out numbers not handled by regex fallback.}
The LLM prompt explicitly says ``eight hundred $\rightarrow$ 800''. But the regex
fallback only matches digit sequences (\texttt{[\textbackslash d,]+}). If LLM is offline,
\textit{"sugar under eight hundred"} gets no price filter. The system silently degrades
without warning.
\end{issueBox}

\begin{issueBox}
\textbf{ISSUE: PKR currency not recognized.}
\textit{"sugar under 50000 PKR"} --- the regex strips the number but does not convert
PKR to USD. The price filter would apply 50000 as a USD ceiling, which is wildly wrong
for a USD/MT comparison (50000 PKR $\approx$ \$180 USD/MT --- reasonable --- but the
comparison is still semantically wrong since \texttt{usd\_per\_mt} is USD).
\end{issueBox}

\begin{warnBox}
\textbf{WARNING: ``bulk'' and ``premium'' not in regex ranking words list.}
As noted above, the regex fallback's \texttt{price\_ranking\_words} set is missing
\textit{``bulk''}, \textit{``premium''}, \textit{``reliable''}, \textit{``top supplier''}, \textit{``established''}.
These are only handled by the LLM. If LLM is offline (missing API key), these signals
produce no \texttt{ranking\_hint} and the ranking uses default weights.
\end{warnBox}

% ========================================================================
\section{Stage 2: Product Resolution --- Deep Analysis}
% ========================================================================

\subsection{The Four-Pass Fuzzy Matching System}

The product resolution is implemented in \texttt{search\_service.py:249} as a sequential
waterfall. Each pass only executes if all prior passes returned 0 matches:

\begin{table}[H]
\centering
\renewcommand{\arraystretch}{1.3}
\begin{tabularx}{\textwidth}{clXc}
\toprule
\textbf{Pass} & \textbf{Method} & \textbf{Examples that Match} & \textbf{Threshold} \\
\midrule
1 & \texttt{istartswith(keyword)} & \textit{"dex"} $\to$ Dextrose; \textit{"sug"} $\to$ Sugar & Exact prefix \\
2 & \texttt{istartswith(keyword)} on items & Same but on ProductItem names & Exact prefix \\
2.5 & \texttt{icontains(keyword)} & \textit{"palm oil"} $\to$ Crude Palm Oil & Substring \\
3 & PostgreSQL Trigram & \textit{"dextos"} $\to$ Dextrose (similarity 0.43) & 0.3 \\
4 & Word-by-word & \textit{"unga bunga dex"} tries each word & per-word trgm \\
\bottomrule
\end{tabularx}
\caption{Product resolution passes and their matching strategies}
\end{table}

\subsection{What Works Well}

\begin{okBox}
Pass 1 handles the dominant pattern of partial abbreviation search efficiently. \\
Pass 2.5 correctly handles infix patterns like \textit{"palm oil"} finding \textit{"Crude Palm Oil"}. \\
Pass 3 (trigram) handles typos with high reliability (PostgreSQL pg\_trgm is accurate). \\
Pass 4 (word-by-word) now uses DB-scored similarity to pick the most product-like word
(replacing the broken length-based heuristic).
\end{okBox}

\subsection{Product Resolution Failures}

\begin{issueBox}
\textbf{CRITICAL: Full-keyword search before decomposition for multi-word qualifiers.}
For keyword \textit{"refined sugar"}: \\
Pass 1: \texttt{istartswith("refined sugar")} --- fails (no product starts with "refined sugar") \\
Pass 2: same --- fails \\
Pass 2.5: \texttt{icontains("refined sugar")} --- fails unless ``Refined Sugar'' exists verbatim in DB \\
Pass 3: Trigram(\textit{"refined sugar"}) vs DB entries --- may find \textit{"Sugar"} at ~0.4 \\
The system relies on trigram for any multi-word query where the qualifier is not in the DB name.
Trigram similarity of 0.3 is very lenient; correctness depends on what DB entries exist.
If the DB has both \textit{"Sugar"} and \textit{"Sugarcane Juice"}, trigram might rank them
identically and return both --- merging unrelated results.
\end{issueBox}

\begin{issueBox}
\textbf{ISSUE: Pass 4 only triggers on queries with a space.}
Line 365: \texttt{if not sc\_qs and not item\_qs and " " in product\_keyword:}. \\
If the garbled keyword has no space (e.g.\ \textit{"dextrozz"}), Pass 4 never runs.
The query falls off the end and returns 0 results even though trigram already ran.
However, Pass 3 already ran trigram on the full keyword --- so \textit{"dextrozz"} would
likely match at similarity $\approx$0.38. This is acceptable.
\end{issueBox}

\begin{issueBox}
\textbf{ISSUE: HS code query not handled in main flow.}
If product\_keyword is \textit{"1702.3090"}, all passes try to match this string against
product \textit{names} in the DB. No pass attempts to match HS codes. A user who types
an HS code directly gets 0 results.
\end{issueBox}

\begin{warnBox}
\textbf{WARNING: Disambiguation threshold of 20 is never explained to the user.}
When $>$20 variants match (rare but possible), the UI shows a product picker. But the
threshold of 20 is hard-coded and arbitrary. There is no rationale in comments. A
search for \textit{"oil"} could match 30+ products (palm oil, canola oil, motor oil,
etc.) and trigger the picker --- good. But \textit{"sugar"} might match only 5 variants
(granulated, refined, raw, brown, powdered) and be merged --- potentially mixing distinct
results.
\end{warnBox}

\subsection{HS Code Variant Disambiguation}

When HS code is provided (from autocomplete click), shared HS codes (e.g.\ 1702.3 shared
by Dextrose Ball, Dextrose Anhydrous, Dextrose Monohydrate) are resolved by:

\begin{enumerate}
  \item Exact \texttt{variant\_name} match (highest priority --- from user click)
  \item Unique-word match (keyword contains a word that appears in only one variant name)
  \item Safe fallback: return ALL variants (merge results)
\end{enumerate}

\begin{okBox}
This is the correct design: merging is better than guessing wrong. Users who need
exact variant isolation click the product picker.
\end{okBox}

\begin{warnBox}
\textbf{Minor issue:} Unique-word match requires keyword length $\geq$ 5 characters.
\textit{"dex"} (3 chars) would always fall through to merge all dextrose variants.
This is intentional but means 3-character prefix queries always see merged results.
\end{warnBox}

% ========================================================================
\section{Stage 3: Worldwide vs.\ Pakistan Scope --- Deep Analysis}
% ========================================================================

\subsection{What the Scope Toggle Actually Does}

The scope parameter is the most architecturally significant parameter in the entire
search system. It determines \textit{which subset of the Transaction table} is queried.
The Transaction table contains two fundamentally separate datasets stored together:

\begin{table}[H]
\centering
\renewcommand{\arraystretch}{1.4}
\begin{tabularx}{\textwidth}{lXXl}
\toprule
\textbf{trade\_type} & \textbf{Direction} & \textbf{Seller is\ldots} & \textbf{Buyer is\ldots} \\
\midrule
IMPORT & Foreign $\to$ Pakistan & Foreign exporter & Pakistani importer \\
EXPORT & Pakistan $\to$ Foreign & Pakistani exporter & Foreign importer \\
\bottomrule
\end{tabularx}
\caption{Transaction table trade types and entity semantics}
\end{table}

The (Intent, Scope) $\to$ (trade\_type, company\_field) mapping is:

\begin{table}[H]
\centering
\renewcommand{\arraystretch}{1.4}
\begin{tabularx}{\textwidth}{llllX}
\toprule
\textbf{Intent} & \textbf{Scope} & \textbf{trade\_type} & \textbf{Shown As} & \textbf{Meaning} \\
\midrule
BUY & WORLDWIDE & IMPORT, excl.\ Pak origin & Seller & I am a Pakistani buyer wanting foreign suppliers \\
BUY & PAKISTAN & EXPORT, Pak origin & Seller & I want a Pakistani supplier (local) \\
SELL & WORLDWIDE & EXPORT, excl.\ Pak dest & Buyer & I am a Pakistani seller wanting foreign buyers \\
SELL & PAKISTAN & IMPORT, dest=Pakistan & Buyer & I want to sell to Pakistani buyers \\
\bottomrule
\end{tabularx}
\caption{Scope-intent-to-query mapping}
\end{table}

\subsection{Why the Scope Toggle Matters}

The toggle is not cosmetic. \textbf{It switches between two completely non-overlapping
datasets.} IMPORT and EXPORT records are separate transaction series with different
origin/destination, different seller/buyer sets, and potentially different price ranges.

Example for ``Sugar (BUY intent)'':
\begin{itemize}
  \item \textbf{Worldwide}: Shows Chinese, Indian, Brazilian sugar exporters (IMPORT records).
        Origin countries: China, India, Brazil. Prices: \$300--\$450/MT.
  \item \textbf{Pakistan}: Shows Pakistani sugar mills (EXPORT records, origin=Pakistan).
        Companies: Shakarganj, JDW Group, etc. Prices: \$350--\$420/MT.
\end{itemize}
These are \textit{entirely different supplier lists}. Without the toggle, you cannot access
the Pakistani domestic supplier dataset.

\subsection{Problems with the Current Scope Implementation}

\begin{issueBox}
\textbf{ISSUE: Scope is not inferred from query context.}
If a user on the Worldwide scope types \textit{"Pakistan sugar supplier"}, the query
mentions Pakistan explicitly. But the country extractor would extract ``Pakistan'' as
a country filter, not as a scope switch. The system would filter the IMPORT dataset
for Pakistan as origin --- but IMPORT records have non-Pakistani origins by definition.
The country filter would return 0 results. The correct behavior would be to switch
to PAKISTAN scope automatically.
\end{issueBox}

\begin{issueBox}
\textbf{ISSUE: Scope UX is opaque to new users.}
The toggle labels are \textit{"Worldwide"} and \textit{"Pakistan"} with no explanation.
A new user has no idea that ``Worldwide'' means foreign suppliers and ``Pakistan'' means
local suppliers. There is no tooltip, help text, or description anywhere in the UI.
Many users likely leave it at the default ``Worldwide'' and never discover local supplier data.
\end{issueBox}

\begin{issueBox}
\textbf{ISSUE: Autocomplete is scope-agnostic.}
The autocomplete endpoint (\texttt{autocomplete\_products}) ranks products by
\texttt{total\_volume} across ALL transactions (both IMPORT and EXPORT, all scopes).
A product like ``Basmati Rice'' might appear prominently in autocomplete with 50,000 MT
total volume, but if the user has Worldwide scope + BUY intent selected, there may be
very few foreign basmati rice suppliers in the IMPORT dataset. The user selects it from
autocomplete and gets unexpectedly few results.
\end{issueBox}

\subsection{Impact of Removing the Worldwide/Pakistan Toggle}

\begin{issueBox}
\textbf{Removing the toggle without an alternative would cause data loss.}
The EXPORT dataset (Pakistani sellers/suppliers) would become entirely inaccessible
if the system defaults to WORLDWIDE (IMPORT data). This affects any Pakistani company
looking for: domestic input suppliers, Pakistani partners for joint ventures, or
tracking Pakistani export activity.
\end{issueBox}

\subsection{How to Achieve the Same Functionality Without the Toggle}

There are three viable strategies:

\textbf{Strategy 1 --- Smart Scope Inference (Recommended)}
\begin{enumerate}
  \item If the query contains ``local'', ``Pakistani'', ``Pakistan'', ``domestic'',
        or a Pakistani city name $\rightarrow$ auto-switch to PAKISTAN scope.
  \item If the query contains a foreign country name $\rightarrow$ auto-switch to WORLDWIDE.
  \item Otherwise, default to WORLDWIDE (most common use case for trade intelligence).
\end{enumerate}
This can be implemented in the NLU engine's \texttt{parse()} method before the
OS filter is built.

\textbf{Strategy 2 --- Unified Results with Source Column}
Query both IMPORT and EXPORT simultaneously and return a merged result set with a
``Source'' column: ``Pakistan'' (local) vs.\ country name (foreign). A UI filter chip
(``Show: All / Local / Foreign'') replaces the toggle. This requires changes in the
aggregator to run two queries and merge.

\textbf{Strategy 3 --- Separate Result Sections in UI}
Keep both queries but display them in two collapsible sections: ``Pakistani Suppliers''
and ``International Suppliers'', both on the same page. Users see both without having
to know about the toggle. The toggle is removed from the search bar and becomes
implicit in the results display.

\begin{okBox}
\textbf{Recommendation:} Strategy 1 (smart scope inference) is the lowest-effort,
highest-impact change. It eliminates the confusing toggle while preserving full
data access and making scope selection feel intelligent.
\end{okBox}

% ========================================================================
\section{Stage 4: ORM Aggregation --- Deep Analysis}
% ========================================================================

\subsection{Query Structure and Correctness}

The aggregation in \texttt{aggregation.py:79} is a single grouped Django QuerySet:

\begin{lstlisting}[language=Python]
results = queryset.values(target_field, country_field).annotate(
    total_volume    = Sum('qty_mt'),
    weighted_price_sum = Sum(F('qty_mt') * F('usd_per_mt')),  # for weighted avg
    shipment_count  = Count('tx_reference', distinct=True),
    last_shipment_date = Max('reporting_date'),
    max_shipment_vol   = Max('qty_mt'),
    avg_shipment_vol   = Avg('qty_mt')
).order_by('-total_volume')
\end{lstlisting}

\begin{okBox}
\textbf{Correct design decisions:}
\begin{itemize}
  \item Weighted average price (Volume $\times$ Price / Volume) instead of simple average.
        This prevents one large cheap shipment from artificially lowering a supplier's
        displayed price.
  \item Post-aggregation price filter: filters on computed avg\_price, not individual rows.
        This means a supplier who had one expensive shipment out of 100 cheap ones still
        passes a ceiling filter if their average is below the ceiling.
  \item \texttt{distinct=True} on shipment\_count prevents duplicate counting.
\end{itemize}
\end{okBox}

\subsection{Aggregation Issues}

\begin{issueBox}
\textbf{ISSUE: country grouping causes supplier splitting.}
The aggregation groups by \texttt{(target\_field, country\_field)}. If a supplier shipped
from both China and India in different transactions (re-export scenarios), they appear as
\textit{two separate entries}: one for China, one for India. The user sees the same company
name twice with different countries. This is confusing in the results list and inflates
the supplier count.
\end{issueBox}

\begin{issueBox}
\textbf{ISSUE: ``Unknown'' origin\_country entries contaminate results.}
The Transaction model has \texttt{default="Unknown"} for origin\_country. Suppliers with
bad data appear as country=``Unknown'' in results, which is confusing. Furthermore, the
aggregation groups (supplier, "Unknown") separately from (supplier, "China"), potentially
splitting the same supplier into two entries.
\end{issueBox}

\begin{warnBox}
\textbf{WARNING: avg\_price calculation bug in aggregation.py.}
\begin{lstlisting}[language=Python]
avg_price = round(wps / tv, 2) if tv > 0 else 0.0
\end{lstlisting}
\texttt{wps} = weighted price sum (sum of qty * price). \texttt{tv} = total volume. This
correctly computes the weighted average. BUT: if \texttt{usd\_per\_mt} is NULL for some
transactions, those rows contribute 0 to \texttt{wps} but their full qty to \texttt{tv},
making the computed avg\_price artificially low. There is no NULL check in the
annotation expression.
\end{warnBox}

\begin{warnBox}
\textbf{WARNING: Volume filter has a broken soft-exclusion condition.}
In aggregation.py:135:
\begin{lstlisting}[language=Python]
if mss < 0.3 * V and total < 0.5 * V:
    continue  # exclude
\end{lstlisting}
A supplier with 1 large shipment of 1000 MT gets \texttt{mss=1000}. If requested
volume V=500, \texttt{mss > 0.3*V}, so they are NOT excluded. But a company with
100 shipments of 1 MT each also has \texttt{mss=1}, which IS excluded (\texttt{mss < 0.3*V}
AND \texttt{total < 0.5*V} if total $<$ 250). The exclusion penalizes frequent small
shippers even if their capacity could aggregate to handle V.
\end{warnBox}

% ========================================================================
\section{Stage 4: Ranking Model --- Deep Analysis}
% ========================================================================

\subsection{What the Ranking Actually Does}

The ranking in \texttt{search\_service.py:617} computes a composite score using
min-max normalization across the current result set:

\begin{equation}
\text{score}_i = w_1 \textbullet{} \frac{v_i - v_{min}}{v_{max} - v_{min}}
               + w_2 \textbullet{} \frac{c_i - c_{min}}{c_{max} - c_{min}}
               + w_3 \textbullet{} \frac{p^{-1}_i - p^{-1}_{min}}{p^{-1}_{max} - p^{-1}_{min}}
               + w_4 \textbullet{} \frac{r_i - r_{min}}{r_{max} - r_{min}}
\end{equation}

Where $v$ = total\_volume, $c$ = shipment\_count, $p$ = avg\_price (inverted, lower = better),
$r$ = last\_shipment\_date (as ordinal integer), and weights depend on the ranking preset:

\begin{table}[H]
\centering
\renewcommand{\arraystretch}{1.3}
\begin{tabular}{lccccl}
\toprule
\textbf{Preset} & $w_1$ (vol) & $w_2$ (cnt) & $w_3$ (price) & $w_4$ (rec) & \textbf{Triggered by} \\
\midrule
price\_asc & 0.2 & 0.2 & 0.5 & 0.1 & ``cheap'', ``under \$X'', ranking\_hint=price\_asc \\
bulk/premium & 0.5 & 0.3 & 0.1 & 0.1 & ``bulk'', ``reliable'', ranking\_hint=price\_desc \\
default & 0.4 & 0.3 & 0.2 & 0.1 & All other queries \\
\bottomrule
\end{tabular}
\caption{Ranking weight presets}
\end{table}

\subsection{The Dead LTR Model}

\begin{issueBox}
\textbf{CRITICAL: The LightGBM LTR model is dead code.}
The project memory and documentation state a ``70\% heuristic + 30\% LightGBM LambdaRank''
ranking ensemble. This is \textbf{false}. The files implementing LTR were deleted:
\begin{itemize}
  \item \texttt{backend/search/services/ranking.py} --- DELETED (shown in git status as \texttt{D})
  \item \texttt{backend/search/services/ranking\_ltr.py} --- DELETED
\end{itemize}
The model file \texttt{backend/search/models/lgbm\_ltr.txt} may still exist, but no
code loads or applies it. The current ranking is 100\% heuristic with no learned signal.
This is a significant gap given the effort of training the LTR model.
\end{issueBox}

\subsection{Ranking Logic Weaknesses}

\begin{issueBox}
\textbf{ISSUE: Recency signal is negligible.}
$w_4 = 0.1$ in all presets. A supplier whose last shipment was in 2020 vs.\ 2023 differs
by $\approx$1095 ordinal days. After min-max normalization, if the range is 2018--2024
($\approx$2190 days), the 2020 supplier gets $\frac{1095}{2190} \approx 0.5$ recency
score while 2023 gets $\frac{1825}{2190} \approx 0.83$. Difference: 0.33. Weighted:
$0.1 \times 0.33 = 0.033$. For a supplier with high volume/count, this 0.033 difference
in final score is dominated by their volume advantage. Stale suppliers appear at the
top of results.
\end{issueBox}

\begin{issueBox}
\textbf{ISSUE: Min-max normalization is relative to the current result set.}
If only 2 suppliers match a query, the lower one gets score~0 and the higher gets
score~1, regardless of their absolute quality. A single supplier who shipped 1 MT gets
vol\_norm=1.0 (best possible) because there are no competitors. This makes scores
incomparable across queries.
\end{issueBox}

\begin{warnBox}
\textbf{WARNING: Price component excluded when avg\_price is 0.}
Code (line 683): \texttt{if r.get("avg\_price"): score += w3 * p\_norm}. \\
Suppliers with no price data skip the price component entirely. In price\_asc mode
($w_3 = 0.5$), these suppliers lose up to 50\% of potential score. This may correctly
deprioritize them, but it also means price-seeking queries unfairly rank data-poor
suppliers last even if they are real and potentially cheap.
\end{warnBox}

\begin{issueBox}
\textbf{ISSUE: No diversity or deduplication in ranking.}
The ranking does not attempt to surface diverse countries or avoid showing the same
company multiple times (see country-grouping split issue in aggregation). If a Chinese
supplier has 10 entries across 10 different Chinese provinces/ports, all 10 could rank
in the top 10, monopolizing the results.
\end{issueBox}

\begin{warnBox}
\textbf{WARNING: ranking\_hint signal from ``bulk'' can conflict with query score preset.}
If the LLM returns \texttt{ranking\_hint=price\_desc} for a \textit{"bulk dextrose"} query,
the preset switches to \textit{bulk/premium} ($w_1=0.5$). But the signal detection
in search\_service.py:622 also checks the raw query for keywords. Since \textit{"bulk"}
is in the keyword list, this correctly triggers bulk/premium. However, if the LLM
\textit{also} detects \texttt{operator=lte} (e.g.\ user says \textit{"bulk dextrose under \$700"}),
the price\_asc preset takes priority (it checks ranking\_hint==price\_asc OR price keywords
first). The conflict means \textit{"bulk under \$700"} always resolves to price\_asc,
ignoring the \textit{"bulk"} volume preference.
\end{warnBox}

% ========================================================================
\section{Caching Architecture}
% ========================================================================

\subsection{Three-Level Cache Design}

\begin{table}[H]
\centering
\renewcommand{\arraystretch}{1.3}
\begin{tabularx}{\textwidth}{lXll}
\toprule
\textbf{Level} & \textbf{Key Composition} & \textbf{TTL} & \textbf{Savings} \\
\midrule
NLU & MD5(query + ui\_context) & 24h & Skips SetFit + LLM call (~2s) \\
Availability & MD5(subcat\_ids + trade\_type + country) & 1h & Skips DB existence check \\
Aggregation & MD5(subcat\_ids + intent + scope + filters) & 1h & Skips heavy GROUP BY query \\
\bottomrule
\end{tabularx}
\caption{Cache levels and their impact}
\end{table}

\begin{okBox}
The three-level design is well-thought-out. The 24h NLU cache is particularly important
since NLU parsing is the bottleneck (SetFit + LLM = 0.5--3s).
\end{okBox}

\subsection{Caching Issues}

\begin{issueBox}
\textbf{ISSUE: No cache invalidation on data ingestion.}
When new transactions are ingested into the DB, the aggregation and availability caches
do not update for up to 1 hour. Fresh trade data appears stale to users.
There is no \texttt{cache.delete()} call anywhere in the ingestion pipeline.
\end{issueBox}

\begin{warnBox}
\textbf{WARNING: In-memory cache loses all data on server restart.}
If Redis is not configured (\texttt{REDIS\_URL} not set), Django falls back to in-memory
caching which is per-process and non-persistent. Every restart means cold-cache for
all users.
\end{warnBox}

\begin{warnBox}
\textbf{WARNING: NLU cache key ignores \texttt{hs\_code} and \texttt{subcat\_id} parameters.}
The NLU cache key is \texttt{f"nlu:\{query\}:\{ui\_context\}"}. Two different requests
with the same query but different \texttt{hs\_code} values share the same NLU cache.
This is correct (NLU result doesn't depend on hs\_code), but it means the cached
NLU result is reused in supplier\_detail and compare views where the context might differ.
\end{warnBox}

% ========================================================================
\section{Frontend: Search UI --- Deep Analysis}
% ========================================================================

\subsection{SearchHome.js Analysis}

\subsubsection{What Works Well}

\begin{okBox}
\begin{itemize}
  \item 250ms debounce on autocomplete prevents excessive API calls.
  \item Product disambiguation banner shows selected product name with a clear X button.
  \item Scope toggle (Worldwide/Pakistan) is prominent and accessible.
  \item Intent pills (\textit{``I want to buy''}, \textit{``Find suppliers''}) pre-fill the search bar.
  \item Outside-click detection correctly closes the autocomplete dropdown.
\end{itemize}
\end{okBox}

\subsubsection{SearchHome Issues}

\begin{issueBox}
\textbf{ISSUE: Typing a new query after selecting from autocomplete clears selectedHsCode.}
In the \texttt{useEffect} on query change (line 43):
\begin{lstlisting}[language=Java]
setSelectedHsCode(null);
setSelectedProductName(null);
\end{lstlisting}
This runs on \textit{every} keystroke. So if a user selects ``Dextrose Anhydrous'' from
autocomplete (which sets selectedHsCode=``1702.3''), then accidentally presses backspace
and re-types the last letter, the hs\_code is lost. The subsequent search runs without
hs\_code, returning a broader result set.
\end{issueBox}

\begin{issueBox}
\textbf{ISSUE: Scope toggle is not persisted across navigation.}
The scope state lives in \texttt{SearchHome} component. When the user submits a search,
the scope is passed as a URL parameter. However, the scope does NOT appear in the
autocomplete API call:
\begin{lstlisting}[language=Java]
await fetch(`\${API_BASE}/api/search/autocomplete/?q=\${q}`)
// Note: scope is NOT included in autocomplete request
\end{lstlisting}
The autocomplete backend ignores scope and ranks by total volume across all scopes.
A Pakistan-scope user sees the same autocomplete suggestions as a Worldwide-scope user.
\end{issueBox}

\begin{warnBox}
\textbf{WARNING: Example queries in the UI are misleading.}
The bottom of SearchHome.js shows example queries:
\begin{itemize}
  \item \textit{``Dextrose suppliers in Pakistan''} --- the system would detect country=Pakistan but
        the scope might be WORLDWIDE. The country filter would apply to IMPORT origin\_country=Pakistan
        which returns 0 results (IMPORT records have non-Pakistani origins).
  \item \textit{``Buy Urea 46\%''} --- the \textit{``46\%''} part is ignored (no handling of
        product grades/specifications in the NLU). Returns Urea results without grade filtering.
\end{itemize}
\end{warnBox}

\subsection{SearchResults.js Analysis}

\subsubsection{What Works Well}

\begin{okBox}
\begin{itemize}
  \item Disambiguation panel cleanly handles multi-variant products.
  \item Comparison selection (up to 4 suppliers) with sticky compare bar is well implemented.
  \item Broad search warning for empty/too-broad queries is helpful.
  \item Sort controls (relevance / price / volume) give user control over ordering.
  \item Left sidebar filters (country, price range, volume) are clean and functional.
\end{itemize}
\end{okBox}

\subsubsection{SearchResults Critical Issues}

\begin{issueBox}
\textbf{BUG: fetchResults uses \texttt{initialQuery} not \texttt{query} state.}
In SearchResults.js:78:
\begin{lstlisting}[language=Java]
const data = await searchService.search(initialQuery, filters);
\end{lstlisting}
\texttt{initialQuery} is derived from URL params at component mount and is a \texttt{const}.
It never changes when the user types in the results-page search bar. The \texttt{query}
state variable updates on typing, but the actual API call always uses the original
URL query. The user appears to be searching but the results don't update until they
press Enter (which navigates to a new URL and remounts the component). This causes
confusion: the user types, the bar shows new text, but results stay the same.
\end{issueBox}

\begin{issueBox}
\textbf{BUG: Client-side price filter conflicts with NLU price filter.}
When a user searches \textit{``sugar under \$700''}, the backend filters to
avg\_price $\leq$ 700. The results page sidebar also has a price filter. If the user
sets sidebar priceMin=\$500, it further narrows the 0--700 range. Fine.
But if the user then clears the query and searches \textit{``sugar''} (no price filter)
and the sidebar priceMax=\$700 is still set from a previous search, the sidebar now
acts as the ONLY price filter. The priceMax state \textit{persists across searches}
because \texttt{clearFilters()} is only called in \texttt{handleSearch}, not on component
re-render with new URL. The user gets different result counts for the same query
depending on what they searched before.

Actually, looking at the code: \texttt{priceMin}, \texttt{priceMax} and \texttt{volumeMin}
are reset in \texttt{handleSearch()} via direct state setters. But on the initial render
they are read from URL params (\texttt{queryParams.get('price\_min')}). Since a normal
search does not include price\_min in the URL, these initialize to \texttt{''} and remain
that way. The persistence bug would only occur if navigating directly to a URL with
price\_min/max params and then changing the search. Low probability but real.
\end{issueBox}

\begin{issueBox}
\textbf{BUG: \texttt{selectedCountry} not reset between searches.}
\texttt{handleSearch()} sets \texttt{selectedCountry(null)} but \texttt{fetchResults}
depends on \texttt{[initialQuery, scope, hsCode, selectedCountry]}. When
\texttt{selectedCountry} changes, \texttt{fetchResults} re-runs with the previous
\texttt{initialQuery}. If the user picks a country filter, then types a new query and
presses Enter, the navigation to new URL remounts with no country. Correct.
But if the user picks a country and then clicks a variant in the disambiguation panel,
\texttt{handleVariantPick} preserves \texttt{selectedCountry} correctly via
\texttt{queryParamsObj.country}. This part is OK.
\end{issueBox}

\begin{issueBox}
\textbf{ISSUE: ``Last active'' date displayed as raw ISO string.}
The \texttt{last\_shipment\_date} field is displayed directly:
\begin{lstlisting}[language=Java]
<span>Last active: {supplier.last_shipment_date}</span>
\end{lstlisting}
This renders as \texttt{2023-12-15} (ISO format). No human-readable formatting
(\textit{``Dec 2023''} or \textit{``3 months ago''}). Minor UX issue.
\end{issueBox}

\begin{issueBox}
\textbf{ISSUE: Alert for comparison limit is a blocking native dialog.}
\begin{lstlisting}[language=Java]
if (prev.length >= 4) {
    alert("You can only compare up to 4 suppliers at once.");
    return prev;
}
\end{lstlisting}
Native \texttt{alert()} is disruptive and breaks the UX flow. Should be an inline toast.
\end{issueBox}

\begin{warnBox}
\textbf{WARNING: Market Snapshot's ``top\_country'' is simply the first result's country.}
In views.py:108:
\begin{lstlisting}[language=Python]
"top_country": mapped_results[0]["country"] if mapped_results else "N/A"
\end{lstlisting}
This is not the top country by volume or frequency --- it is literally the country of
the \#1 ranked supplier. If the top supplier happens to be from a small country, the
market snapshot shows that country as ``top origin'' even though most volume comes
from elsewhere. This is misleading.
\end{warnBox}

% ========================================================================
\section{Supplier Detail Page --- Deep Analysis}
% ========================================================================

\subsection{hs\_code Bug in Supplier Detail and Compare Views}

\begin{issueBox}
\textbf{BUG: \texttt{hs\_code} URL parameter ignored in supplier-detail and compare.}
In \texttt{views.py:185--190} (supplier\_detail view):
\begin{lstlisting}[language=Python]
subcat_ids, _, product_item_ids = self.search_service._resolve_subcategories(
    product_keyword=parsed_query.get("product", ""),
    hs_code=parsed_query.get("hs_code", ""),  # BUG
    ...
)
\end{lstlisting}
\texttt{parsed\_query} is the NLU result dict. Its keys are: \textit{intent, product,
product\_keyword, country, quantity, price\_filter, os\_filter, entities, ui\_context}.
There is \textbf{no ``hs\_code'' key} in the NLU result. This always passes \texttt{hs\_code=""}
to resolve\_subcategories, bypassing HS-code-based variant resolution entirely.
The correct code should use \texttt{request.query\_params.get('hs\_code', '')}.
The same bug exists in the compare view (line 241).
\end{issueBox}

\subsection{Intelligence Metric Issues}

\begin{warnBox}
\textbf{WARNING: Momentum label uses fixed 90/180 day windows from \textit{today}.}
The momentum calculation compares volume in the last 90 days vs.\ the 90 days before that.
If the dataset ends in 2023 and today is 2026, \texttt{vol\_recent} = 0 and
\texttt{vol\_prev} = 0 for all suppliers. Every supplier gets \texttt{momentum\_label = "Declining"}
because \texttt{vol\_recent == 0 and total\_vol > 0} (line 502--503). This makes the
intelligence metrics meaningless for any dataset that is more than 6 months old.
\end{warnBox}

\begin{warnBox}
\textbf{WARNING: Pricing label uses only the first and last transaction.}
\begin{lstlisting}[language=Python]
price_change = (latest_tx.usd_per_mt - earliest_tx.usd_per_mt) / earliest_tx.usd_per_mt
\end{lstlisting}
If the earliest transaction happened to have an outlier price (e.g.\ a distress sale
at \$200 when normal price is \$700), the computed \texttt{price\_change} would be
enormous and misleading. Using two individual transactions instead of moving averages
makes this metric highly noise-sensitive.
\end{warnBox}

% ========================================================================
\section{Query Families --- Coverage Assessment}
% ========================================================================

The following table assesses all realistic query patterns against the actual system behavior:

\begin{longtable}{lp{5cm}ccc}
\toprule
\textbf{\#} & \textbf{Query Pattern} & \textbf{Intent} & \textbf{Product} & \textbf{Country} \\
\midrule
\endhead
1 & \textit{"dextrose anhydrous"} (bare noun) & \cellOK{BUY} & \cellOK{OK} & \cellOK{N/A} \\
2 & \textit{"buy dextrose"} & \cellOK{BUY} & \cellOK{OK} & \cellOK{N/A} \\
3 & \textit{"sell wheat to UAE"} & \cellOK{SELL} & \cellOK{OK} & \cellOK{UAE} \\
4 & \textit{"find sellers of basmati rice"} & \cellOK{BUY} & \cellOK{OK} & \cellOK{N/A} \\
5 & \textit{"find buyers for rice"} & \cellOK{SELL} & \cellOK{OK} & \cellOK{N/A} \\
6 & \textit{"cheap sugar under \$400"} & \cellOK{BUY} & \cellOK{OK} & \cellOK{N/A} \\
7 & \textit{"bulk wheat from India"} & \cellOK{BUY} & \cellOK{OK} & \cellOK{India} \\
8 & \textit{"import dextrose from china"} & \cellOK{BUY} & \cellOK{OK} & \cellOK{China} \\
9 & \textit{"dextos"} (typo) & \cellOK{BUY} & \cellOK{Trigram} & \cellOK{N/A} \\
10 & \textit{"dex"} (abbreviation) & \cellOK{BUY} & \cellOK{Prefix} & \cellOK{N/A} \\
11 & \textit{"organic wheat"} & \cellOK{BUY} & \cellWarn{Partial} & \cellOK{N/A} \\
12 & \textit{"refined sugar"} & \cellOK{BUY} & \cellWarn{Partial} & \cellOK{N/A} \\
13 & \textit{"sugar and wheat"} & \cellOK{BUY} & \cellBad{Single} & \cellOK{N/A} \\
14 & \textit{"dextrose for sale"} & \cellBad{BUY} & \cellOK{OK} & \cellOK{N/A} \\
15 & \textit{"I am a supplier of rice"} & \cellBad{BUY} & \cellOK{OK} & \cellOK{N/A} \\
16 & \textit{"1702.3090"} (HS code) & \cellOK{BUY} & \cellBad{No match} & \cellOK{N/A} \\
17 & \textit{"SEAWALL"} (company name) & \cellOK{BUY} & \cellBad{No match} & \cellOK{N/A} \\
18 & \textit{"bangladesh wheat"} & \cellOK{BUY} & \cellOK{OK} & \cellBad{Dropped} \\
19 & \textit{"myanmar rice"} & \cellOK{BUY} & \cellOK{OK} & \cellBad{Dropped} \\
20 & \textit{"sugar under eight hundred"} & \cellOK{BUY} & \cellOK{OK} & \cellBad{No filter} \\
21 & \textit{"sugar 50000 PKR"} & \cellOK{BUY} & \cellOK{OK} & \cellBad{USD mismatch} \\
22 & \textit{"who needs sugar"} & \cellOK{BUY} & \cellOK{OK} & \cellOK{N/A} \\
23 & \textit{"sugar available to export"} & \cellOK{SELL} & \cellOK{OK} & \cellOK{N/A} \\
24 & \textit{"top 5 dextrose suppliers"} & \cellOK{BUY} & \cellOK{OK} & \cellOK{top\_k=5} \\
25 & \textit{Urdu script query (e.g.\ buy sugar in Urdu)} & \cellWarn{BUY} & \cellBad{No match} & \cellBad{None} \\
\bottomrule
\caption{Query family coverage matrix (Partial = fragile/trigram-dependent; Dropped = silently ignored)}
\end{longtable}

\textbf{Summary:} 16/25 queries fully correct; 5/25 partially handled; 4/25 broken.
Effective correct rate: approximately 64\%.

% ========================================================================
\section{Database and Indexing}
% ========================================================================

\subsection{Index Analysis}

\begin{table}[H]
\centering
\renewcommand{\arraystretch}{1.3}
\begin{tabularx}{\textwidth}{Xll}
\toprule
\textbf{Index} & \textbf{Query it serves} & \textbf{Adequacy} \\
\midrule
\texttt{tx\_type\_product\_idx} (trade\_type, product\_item) & Main aggregation WHERE clause & \cellOK{Excellent} \\
\texttt{tx\_prod\_type\_seller\_idx} (product\_item, trade\_type, seller) & BUY aggregation GROUP BY & \cellOK{Excellent} \\
\texttt{tx\_prod\_type\_buyer\_idx} (product\_item, trade\_type, buyer) & SELL aggregation GROUP BY & \cellOK{Excellent} \\
\texttt{reporting\_date} & Time filter & \cellOK{Good} \\
\texttt{origin\_country} & Country filter (import) & \cellOK{Good} \\
\texttt{destination\_country} & Country filter (export) & \cellOK{Good} \\
\texttt{seller} & supplier-detail filter & \cellWarn{No lower case} \\
\texttt{buyer} & buyer-detail filter & \cellWarn{No lower case} \\
\midrule
\textit{Missing: ProductSubCategory.name} & Prefix/contains search & \cellBad{No index} \\
\textit{Missing: ProductSubCategory.hs\_code} & HS code lookup & \cellBad{No index} \\
\textit{Missing: pg\_trgm GIN on name} & Trigram search & \cellBad{Sequential scan} \\
\bottomrule
\end{tabularx}
\caption{Database index adequacy assessment}
\end{table}

\begin{issueBox}
\textbf{CRITICAL: No GIN trigram index on ProductSubCategory.name.}
Pass 3 of product resolution uses \texttt{TrigramSimilarity('name', keyword)} which
requires a GIN index on the column with the \texttt{pg\_trgm} extension for performance.
Without it, every trigram search does a sequential scan of all product names. For a
small product catalog this is acceptable, but as the catalog grows this will become
a bottleneck. The fix is:
\begin{lstlisting}[language=SQL]
CREATE INDEX idx_subcat_name_trgm ON trade_data_productsubcategory
USING GIN (name gin_trgm_ops);
\end{lstlisting}
\end{issueBox}

\begin{warnBox}
\textbf{WARNING: seller/buyer indexes cannot accelerate \texttt{\_\_iexact} lookups.}
The supplier-detail query uses \texttt{seller\_\_iexact=seller\_name.strip()}.
The standard B-tree index on \texttt{seller} is case-sensitive. An \texttt{iexact} lookup
performs a sequential scan with a \texttt{LOWER()} function call, not an index seek.
Fix: create a functional index on \texttt{LOWER(seller)}.
\end{warnBox}

% ========================================================================
\section{The Worldwide/Pakistan Toggle: Removal Analysis}
% ========================================================================

\subsection{Full Technical Impact Assessment}

To remove the Worldwide/Pakistan toggle and achieve equivalent functionality, the following
changes would be needed across the stack:

\begin{table}[H]
\centering
\renewcommand{\arraystretch}{1.3}
\begin{tabularx}{\textwidth}{lXl}
\toprule
\textbf{Layer} & \textbf{Change Required} & \textbf{Complexity} \\
\midrule
NLU Engine & Add scope inference from query signals (\textit{``local''}, \textit{``Pakistani''}, foreign country) & Medium \\
Search Service & Pass inferred scope to aggregation instead of UI parameter & Low \\
Aggregation & Optionally run both queries and merge (Strategy 2) & High \\
Views & Remove \texttt{scope} URL parameter; add inferred scope to response & Low \\
Frontend & Remove scope toggle buttons & Trivial \\
Frontend & Handle cases where both scopes return results & Medium \\
Autocomplete & Scope-aware ranking (show by relevant scope volume) & Medium \\
\bottomrule
\end{tabularx}
\caption{Changes required to remove the scope toggle}
\end{table}

\subsection{Recommended Implementation: Smart Scope Inference}

Add a \texttt{\_infer\_scope()} method to the NLU pipeline:

\begin{lstlisting}[language=Python]
def _infer_scope(query: str, resolved_country: str) -> str:
    """Infer WORLDWIDE or PAKISTAN from query signals."""
    q = query.lower()

    local_signals = {"local", "domestic", "pakistan", "pakistani", "karachi",
                     "lahore", "islamabad", "faisalabad", "sialkot"}

    # Explicit local signal in query
    if any(s in q for s in local_signals):
        return "PAKISTAN"

    # A specific foreign country was mentioned $\rightarrow$ Worldwide
    if resolved_country and resolved_country.lower() != "pakistan":
        return "WORLDWIDE"

    # Default: Worldwide (show foreign suppliers/buyers)
    return "WORLDWIDE"
\end{lstlisting}

This single function, integrated into \texttt{nlu\_engine.parse()}, enables removing the
toggle entirely while automatically switching context based on the query's natural language.
The scope would then be part of the NLU result and flow naturally to the aggregator.

\textbf{Trade-offs:}
\begin{itemize}
  \item Users lose explicit control over scope --- power users may prefer the toggle.
  \item A query like \textit{``sugar suppliers''} always goes to Worldwide. Users wanting
        Pakistani suppliers must remember to say ``local''.
  \item Recommendation: keep the toggle as an advanced option (collapsed by default)
        while defaulting to smart inference.
\end{itemize}

% ========================================================================
\section{Consolidated Issue Register}
% ========================================================================

\begin{longtable}{clp{7cm}l}
\toprule
\textbf{\#} & \textbf{Layer} & \textbf{Issue} & \textbf{Severity} \\
\midrule
\endhead
1 & Ranking & LTR model code deleted; ranking is heuristic-only & \cellBad{Critical} \\
2 & Ranking & Recency weight 0.1 makes stale suppliers rank high & \cellBad{High} \\
3 & Ranking & Rel. min-max scores not comparable across queries & \cellBad{High} \\
4 & Intent & ``for sale'', ``in stock'' signals not detected & \cellBad{High} \\
5 & Intent & Urdu/non-English queries not handled & \cellBad{High} \\
6 & Country & 27-country list misses Bangladesh, Philippines, Myanmar, Russia, etc. & \cellBad{High} \\
7 & Product & HS code as query not handled & \cellBad{High} \\
8 & Product & Company name as query returns 0 results & \cellBad{High} \\
9 & Bug & \texttt{hs\_code} ignored in supplier-detail \& compare views & \cellBad{High} \\
10 & Bug & Momentum labels always ``Declining'' for old datasets & \cellBad{High} \\
11 & DB & No GIN trigram index on product names (sequential scan) & \cellBad{High} \\
12 & DB & \texttt{seller/buyer} iexact bypasses B-tree index & \cellWarn{Medium} \\
13 & Aggregation & Country grouping splits single supplier into multiple rows & \cellWarn{Medium} \\
14 & Aggregation & NULL usd\_per\_mt rows deflate weighted avg price & \cellWarn{Medium} \\
15 & Product & Multi-word qualifiers (``organic'', ``refined'') break extraction & \cellWarn{Medium} \\
16 & Product & Multi-product queries unsupported & \cellWarn{Medium} \\
17 & NLU & ``bulk'' ranking hint lost when LLM is offline & \cellWarn{Medium} \\
18 & Scope & ``Pakistan'' in query doesn't auto-switch to Pakistan scope & \cellWarn{Medium} \\
19 & Autocomplete & Scope-agnostic; shows products with 0 results for current scope & \cellWarn{Medium} \\
20 & Frontend & \texttt{initialQuery} vs \texttt{query} state mismatch in results page & \cellWarn{Medium} \\
21 & Frontend & \texttt{alert()} for compare limit is blocking dialog & \cellWarn{Low} \\
22 & Frontend & Last-active date not human-formatted & \cellWarn{Low} \\
23 & Frontend & Scope not persisted to autocomplete API call & \cellWarn{Low} \\
24 & Cache & No cache invalidation on data ingestion & \cellWarn{Medium} \\
25 & Insight & \texttt{top\_country} in market snapshot is first result's country, not top by volume & \cellWarn{Low} \\
26 & Insight & Pricing label based on two transactions only (noise-sensitive) & \cellWarn{Medium} \\
27 & Scope & Pakistan/Worldwide toggle has no UX explanation & \cellWarn{Low} \\
\bottomrule
\caption{Consolidated issue register with severity ratings}
\end{longtable}

% ========================================================================
\section{Recommendations --- Prioritized}
% ========================================================================

\subsection{Priority 1 --- Ranking Restoration (Critical)}

\begin{enumerate}
  \item \textbf{Re-integrate the LightGBM LTR model.} The model file exists at
        \texttt{search/models/lgbm\_ltr.txt}. Re-create \texttt{ranking.py} that loads
        the model and applies it as a re-ranker on the top-100 heuristic results.
        If model quality is uncertain, blend: 60\% heuristic + 40\% LTR.
  \item \textbf{Increase recency weight to $w_4 = 0.25$} in the default preset to
        make stale suppliers rank lower.
  \item \textbf{Cap result set at top 50 by total\_volume before ranking} to focus
        the score computation on relevant results.
\end{enumerate}

\subsection{Priority 2 --- NLU Coverage Gaps}

\begin{enumerate}
  \item Expand STANDARD\_COUNTRIES to at least 80 countries (use ISO 3166-1 list).
  \item Add selling-signal patterns to intent detection: \textit{``for sale''}, \textit{``in stock''},
        \textit{``we have X''}, \textit{``X available''}.
  \item Add HS code detection to \texttt{extract\_product\_keyword}: if query is all digits
        or digit-dot pattern, treat as HS code and bypass NLU to do direct DB lookup.
  \item Add \textit{``bulk''}, \textit{``premium''}, \textit{``reliable''} to the regex
        \texttt{price\_ranking\_words} set as fallback for when LLM is offline.
\end{enumerate}

\subsection{Priority 3 --- Bug Fixes}

\begin{enumerate}
  \item Fix \texttt{hs\_code} in supplier-detail view: replace
        \texttt{parsed\_query.get("hs\_code", "")} with
        \texttt{request.query\_params.get('hs\_code', '')}.
  \item Fix Momentum labels to use relative dates: instead of computing from today,
        compute from the last transaction date in the dataset.
  \item Create GIN trigram index on \texttt{ProductSubCategory.name}.
  \item Create functional index \texttt{LOWER(seller)} and \texttt{LOWER(buyer)}.
\end{enumerate}

\subsection{Priority 4 --- Scope UX Improvements}

\begin{enumerate}
  \item Add smart scope inference to NLU engine (see Section~11.2).
  \item Add tooltip/help text to Worldwide/Pakistan toggle explaining what each means.
  \item Make autocomplete scope-aware (pass scope to autocomplete endpoint).
  \item Replace \texttt{alert()} with inline toast notification in compare flow.
\end{enumerate}

\subsection{Priority 5 --- Long-term Quality}

\begin{enumerate}
  \item Fix country grouping in aggregation: group by company only, aggregate
        country as a list or select the most frequent.
  \item Add cache invalidation hook in the data ingestion pipeline.
  \item Implement basic Urdu transliteration/detection as a preprocessing step.
  \item Add a company name search path that queries the Transaction table directly
        by seller/buyer name when the product keyword matches no products.
  \item Fix Market Snapshot \texttt{top\_country}: compute by
        volume-weighted frequency, not first result's country.
\end{enumerate}

% ========================================================================
\section{Appendix: Data Flow Diagram}
% ========================================================================

\begin{center}
\textbf{Complete Search Request Data Flow}

\vspace{1em}
\begin{tabular}{|l|}
\hline
\textbf{1. HTTP GET /api/search/?q=``cheap dextrose china''\ \&scope=worldwide} \\
\hline
\end{tabular}

$\downarrow$

\begin{tabular}{|l|}
\hline
\textbf{2. NLU Engine (nlu\_engine.py)} \\
Intent: BUY (``buy\textbackslash w*'' regex matches nothing; default BUY) \\
Country: China (FALLBACK\_COUNTRIES dict) \\
Keyword: ``dextrose'' (stop-word strip removes ``cheap'') \\
Price: \{ranking\_hint: price\_asc\} (``cheap'' detected in regex) \\
\hline
\end{tabular}

$\downarrow$

\begin{tabular}{|l|}
\hline
\textbf{3. Scope Mapping} \\
ui\_context=``worldwide'', intent=BUY $\rightarrow$ trade\_type=IMPORT, excl.\ origin=Pakistan \\
\hline
\end{tabular}

$\downarrow$

\begin{tabular}{|l|}
\hline
\textbf{4. Product Resolution (search\_service.py:249)} \\
Pass 1: istartswith(``dextrose'') $\rightarrow$ finds: Dextrose Anhydrous, Dextrose Ball, Dextrose Monohydrate \\
Availability check: filter those with IMPORT records, origin\_country icontains ``China'' \\
Returns: subcat\_ids=[12, 13, 14], variant\_list=[...] \\
\hline
\end{tabular}

$\downarrow$

\begin{tabular}{|l|}
\hline
\textbf{5. ORM Aggregation (aggregation.py:8)} \\
SELECT seller, origin\_country, SUM(qty\_mt), ... \\
FROM transaction WHERE product\_item\_sub\_category\_id IN (12,13,14) \\
AND trade\_type='IMPORT' AND origin\_country != 'Pakistan' \\
AND origin\_country IN ['China'] \\
GROUP BY seller, origin\_country ORDER BY total\_volume DESC \\
\hline
\end{tabular}

$\downarrow$

\begin{tabular}{|l|}
\hline
\textbf{6. Ranking (search\_service.py:617)} \\
Preset: price\_asc (w1=0.2, w2=0.2, w3=0.5, w4=0.1) \\
Min-max normalize: volume, count, 1/price, recency \\
Sort by composite score DESC \\
\hline
\end{tabular}

$\downarrow$

\begin{tabular}{|l|}
\hline
\textbf{7. API Response} \\
\{results: [...], market\_snapshot: \{avg\_price: 680, top\_country: ``China''\}, count: 8\} \\
\hline
\end{tabular}
\end{center}

% ========================================================================
\section{Conclusion}
% ========================================================================

The Zarailink search engine is a well-structured, resilient system that handles the
majority of straightforward trade queries correctly. Its greatest strengths are the
correct aggregation semantics (weighted average price, post-aggregation filtering),
the three-level caching architecture, and the graceful degradation chain (LLM $\rightarrow$ regex,
OpenSearch $\rightarrow$ ORM, SetFit $\rightarrow$ keyword regex).

Its most serious weaknesses are: (1) the deleted LTR model leaving ranking purely
heuristic; (2) significant NLU blind spots on any non-English query, implicit selling
signals, and multi-word product qualifiers; (3) a country coverage list covering only
27 of 195 countries; (4) an opaque scope toggle with no user explanation.

For a production-ready, commercially viable trade intelligence platform, the critical
issues --- particularly the ranking model, the hs\_code bug in supplier-detail, and the
momentum labels using today's date --- should be addressed before the next major user demo.

\vspace{2em}
\begin{center}
\rule{6cm}{0.4pt}\\
\textit{End of Report}\\
\rule{6cm}{0.4pt}
\end{center}

\end{document}
